% !PRINTED 626
\weberpage{626}
b) Ist umgekehrt $\gamma$ durch $q-1$ theilbar, so folgt aus (8), dass $p-1$ durch $q$ theilbar ist. Nehmen wir also $p-1$ nicht durch $q$ theilbar an, so ist auch $e$ nicht durch $q-1$ theilbar, und $\xi$ durchläuft die Reihe der Zahlen (9).

Bilden wir die Summe dieser Zahlen $\xi$, so folgt:
\[ (g-1) \Sigma \xi \equiv g^e - 1 \, (\mathrm{mod}\ m), \]
und diese Zahl ist durch $q$ theilbar oder nicht theilbar, je nachdem $e$ durch $q-1$ theilbar oder nicht theilbar ist. Da ausserdem $g-1$ nicht durch $q$ theilbar ist, so ist $\Sigma\xi$ durch $q$ theilbar oder nicht theilbar, je nachdem $p-1$ durch $q$ theilbar oder nicht theilbar ist.

2. Wir wenden uns zu dem Falle, dass $m = 2^\varkappa$ eine Potenz von 2 ist, und hier haben wir wieder zu unterscheiden, ob $p-1$ durch 4 oder nur durch 2 theilbar ist.

a) Ist $p-1$ durch 4 theilbar, so ist (\S.\,15):
\[ p \equiv 5^\beta \, (\mathrm{mod}\ m), \]
und $f$ ist die kleinste positive Zahl, die der Bedingung
\[ \beta f \equiv 0 \, [\mathrm{mod}\ \tfrac{1}{2}\varphi(m)] \]
genügt. Ist $f=1$, also $\beta$ durch $\tfrac{1}{2}\varphi(m) = 2^{\varkappa-2}$ theilbar, so ist $e = \varphi(m)$, und $\xi$ durchläuft die ganze Gruppe $\mathfrak{N}$, d.\,h. alle ungeraden Zahlen zwischen 0 und $m$. Dies ist der Fall, wo $p-1$ durch $m$ theilbar ist.

Ist $2^{\varkappa_1-2}$ die höchste Potenz von 2, die in $\beta$ aufgeht, und $2 \leqq \varkappa_1 < \varkappa$, so ist, wenn wir $m_1 = 2^{\varkappa_1}$ und $m = m_1 m_2$ setzen,
\[ f = 2^{\varkappa - \varkappa_1}, \quad e = 2^{\varkappa_1 - 1} = \varphi(m_1), \]
und es ist, wenn $s$ eine ungerade Zahl ist,
\[ p - 1 \equiv 5^{s\tfrac{1}{2}\varphi(m_1)} - 1 \ldots \, (\mathrm{mod}\ m). \]
% [sic] Source prints three low dots after -1 in this congruence; no continuation term is specified.
Nach \S.\,15 ist $5^{\tfrac{1}{2}\varphi(m_1)} - 1$ durch $m_1$, aber durch keine höhere Potenz von 2 theilbar. Da man ausserdem nach denselben Sätzen
\[ 5^{s\tfrac{1}{2}\varphi(m_1)} \equiv 5^{\tfrac{1}{2}\varphi(m_1)} \, (\mathrm{mod}\ 2m_1) \]
hat, so folgt, dass $m_1$ der grösste gemeinschaftliche Theiler von $p-1$ und $m$ ist.

Wenn wir also, wie oben, den Körper $\Omega_{m_1}$ zuziehen, so ergiebt sich, dass, wenn $m_1 \geqq 4$ der grösste gemeinschaftliche Theiler von $p-1$ und $m$ ist, $\xi$ ein volles System ungerader Reste von $m_1$ durchläuft, und dass $\mathfrak{p}_\xi$ zugleich die Primfactoren von $p$ im Körper $\Omega_{m_1}$ sind.
